% UET Thesis Pandoc LaTeX Template
% Customized for University of Engineering and Technology, VNU

\clearpage
\phantomsection

\setcounter{chapter}{3}
\chapter[Sơ đồ chữ ký số]{Sơ đồ chữ ký số}

Chương này trình bày ba sơ đồ chữ ký số được xây dựng trên nền tảng
đường cong elliptic đã sinh ở Chương 3: \textbf{Schnorr},
\textbf{ECDSA}, và \textbf{chữ ký tổng hợp BLS}. Với mỗi sơ đồ, chúng ta
đi qua lý thuyết cốt lõi, phân tích cấu trúc dữ liệu và cài đặt thực tế
bằng Rust. Kết quả đo hiệu năng và phân tích bảo mật được trình bày
riêng tại \textbf{Chương 5}.

\section{Cặp khóa (KeyPair)}\label{cux1eb7p-khuxf3a-keypair}

Mọi sơ đồ chữ ký đều dùng chung cấu trúc cặp khóa:

\begin{Shaded}
\begin{Highlighting}[]
\KeywordTok{pub} \KeywordTok{struct}\NormalTok{ KeyPair}\OperatorTok{\textless{}}\NormalTok{C}\OperatorTok{:}\NormalTok{ SWCurveConfig}\OperatorTok{\textgreater{}} \OperatorTok{\{}
    \KeywordTok{pub}\NormalTok{ private\_key}\OperatorTok{:}\NormalTok{ U1024}\OperatorTok{,}        \CommentTok{// d in [1, r{-}1], bí mật}
    \KeywordTok{pub}\NormalTok{ public\_key}\OperatorTok{:}\NormalTok{  AffinePoint}\OperatorTok{\textless{}}\NormalTok{C}\OperatorTok{\textgreater{},} \CommentTok{// Q = [d]G, công khai}
\OperatorTok{\}}
\end{Highlighting}
\end{Shaded}

\captionof{lstlisting}{Cấu trúc dữ liệu cặp khóa KeyPair}\label{lst:ch4_keypair}\vspace{12pt}

\begin{itemize}
\tightlist
\item
  \textbf{Khóa riêng} \texttt{d}: một số nguyên ngẫu nhiên trong
  \([1, r-1]\), sinh bởi CSPRNG (\texttt{U1024::rand}).
\item
  \textbf{Khóa công khai} \texttt{Q\ =\ {[}d{]}G}: tính bằng phép nhân
  vô hướng. Với \(r \approx 2^{512}\), tính \(Q\) từ \(d\) là tính khả
  thi (mili giây), nhưng bài toán ngược (tính \(d\) từ \(Q\)) là ECDLP
  --- độ phức tạp \(O(2^{256})\).
\end{itemize}

Cặp khóa được lưu và đọc từ file nhị phân 128 byte (big-endian) với
quyền truy cập khắt khe \texttt{0600} trên hệ điều hành Unix phân quyền
để giảm rủi ro bị đọc lén bởi các tài khoản người dùng khác trên cùng
máy chủ. Các thao tác khởi tạo (\texttt{generate}) ngẫu nhiên và truy
xuất dữ liệu tĩnh (\texttt{save}) được xử lý bằng API tiêu chuẩn. Chi
tiết mã nguồn quản lý vòng đời cặp khóa được đính kèm tại mã nguồn công
khai của dự án \citep{LumenMath}.

\subsection{Hàm băm mở rộng}\label{huxe0m-bux103m-mux1edf-rux1ed9ng}

Vì \(r \approx 2^{512}\) và SHA-256 chỉ cho 256 bit, hệ thống dùng
\textbf{4 lần SHA-256} với tiền tố khác nhau để sinh 1024 bit, rồi cắt
lấy 512 bit thấp khi cần:

Để đáp ứng yêu cầu kích thước, hàm \texttt{hash\_message} được cấu hình
để thực hiện băm lặp chuỗi đầu vào theo các bộ tiền tố biến thiên. Cấu
trúc lặp này vận hành tương đương như một hàm \textbf{SHA-1024} tùy
chỉnh. Sự thiết kế này đảm bảo vùng không gian đầu ra đủ độ rộng nhằm
phòng tránh hiệu quả các biến thế đụng độ vòng lặp có nghĩa với mức bảo
mật \(r \approx 2^{512}\). Chi tiết cài đặt thuật toán kết hợp băm mở
rộng được cung cấp tại mã nguồn công khai của dự án \citep{LumenMath}.

\section{Chữ ký Schnorr}\label{chux1eef-kuxfd-schnorr}

\subsection{Lý thuyết}\label{luxfd-thuyux1ebft}

Schnorr là sơ đồ chữ ký Sigma-protocol dựa trên giả thuyết logarit rời
rạc (DL) \citep{Schnorr1989}. Nó được ưa chuộng vì độ đơn giản của chứng
minh bảo mật (trong mô hình Oracle ngẫu nhiên) và hỗ trợ tổng hợp tự
nhiên.

\textbf{Ký:} Cho khóa riêng \(d \in [1, r-1]\) và thông điệp \(m\):

\begin{enumerate}
\def\labelenumi{\arabic{enumi}.}
\tightlist
\item
  Chọn ngẫu nhiên \(k \xleftarrow{R} [1, r-1]\)
\item
  Tính điểm \(R = [k]G\)
\item
  Tính thách thức \(e = H(R_x \| R_y \| m) \pmod{r}\)
\item
  Tính \(s = k + e \cdot d \pmod{r}\)
\item
  Chữ ký: \(\sigma = (R, s)\)
\end{enumerate}

\textbf{Xác minh:} Cho khóa công khai \(Q\) và chữ ký \((R, s)\):

\begin{enumerate}
\def\labelenumi{\arabic{enumi}.}
\tightlist
\item
  Tính \(e = H(R_x \| R_y \| m)\)
\item
  Kiểm tra \([s]G = R + [e]Q\)
\end{enumerate}

Tính đúng đắn: \([s]G = [k + ed]G = [k]G + [e][d]G = R + [e]Q\) {[}ok{]}

\subsection{Cấu trúc dữ liệu}\label{cux1ea5u-truxfac-dux1eef-liux1ec7u}

\begin{Shaded}
\begin{Highlighting}[]
\KeywordTok{pub} \KeywordTok{struct}\NormalTok{ SchnorrSignature}\OperatorTok{\textless{}}\NormalTok{C}\OperatorTok{:}\NormalTok{ SWCurveConfig}\OperatorTok{\textgreater{}} \OperatorTok{\{}
    \KeywordTok{pub}\NormalTok{ r\_point}\OperatorTok{:}\NormalTok{ AffinePoint}\OperatorTok{\textless{}}\NormalTok{C}\OperatorTok{\textgreater{},}  \CommentTok{// R = [k]G  (điểm 256 byte: 2x128 byte tọa độ)}
    \KeywordTok{pub}\NormalTok{ s}\OperatorTok{:}\NormalTok{       U1024}\OperatorTok{,}           \CommentTok{// s in [0, r{-}1], 128 byte}
\OperatorTok{\}}
\end{Highlighting}
\end{Shaded}

\captionof{lstlisting}{Cấu trúc dữ liệu chữ ký Schnorr}\label{lst:ch4_schnorr}\vspace{12pt}

Kích thước chữ ký: \(2 \times 128 + 128 = \mathbf{384}\) byte (R gồm hai
tọa độ \(x, y\) mỗi tọa độ 128 byte, cộng thêm \(s\) 128 byte).

\subsection{Cài đặt}\label{cuxe0i-ux111ux1eb7t}

Các bước ký và xác minh thuật toán Schnorr được nhóm tác giả cài đặt bám
sát với quy trình lý thuyết toán học (đã trình bày ở Mục 4.2.1), tận
dụng triệt để kiến trúc cấp phát và an toàn bộ nhớ của Rust để loại bỏ
lỗi tràn giới hạn (buffer overflow) sinh ra từ việc xử lý chuỗi số lớn.
Thuật toán cài đặt tích hợp hàm băm nội sinh \texttt{challenge\_hash}
dùng nối trực tiếp các thành phần \(R_x, R_y, m\) trước khi gọi hàm băm
mở rộng. Quy trình gộp này đảm bảo ``thách thức'' sinh ra phụ thuộc định
danh vào toàn bộ ngữ cảnh thông điệp, triệt tiêu mọi rủi ro tấn công giả
mạo (forgery) bóc tách dữ liệu --- đặc biệt theo kiểu công kích
Pohlig-Hellman trên phần tử nonce.

Chi tiết mã nguồn của các hàm \texttt{sign()} và \texttt{verify()} thuộc
thuật toán Schnorr được đính kèm tại mã nguồn công khai của dự án
\citep{LumenMath}.

\subsection{Bảo mật}\label{bux1ea3o-mux1eadt}

\begin{itemize}
\tightlist
\item
  \textbf{Unforgeability (EUF-CMA):} Chứng minh được trong Random Oracle
  Model dưới giả thuyết DL. Kẻ tấn công cần giải ECDLP hoặc bẻ hàm băm.
\item
  \textbf{Nguy cơ nonce tái sử dụng:} Nếu \(k\) bị tái sử dụng cho hai
  thông điệp khác nhau, kẻ công phá có thể tính
  \(d = (s_1 - s_2)/(e_1 - e_2) \pmod{r}\). Cài đặt dùng CSPRNG
  (\texttt{rand::thread\_rng}) cho mỗi lần ký.
\end{itemize}

\section{Chữ ký ECDSA}\label{chux1eef-kuxfd-ecdsa}

\subsection{Lý thuyết}\label{luxfd-thuyux1ebft-1}

ECDSA (Elliptic Curve Digital Signature Algorithm) là chuẩn IEEE P1363
và FIPS 186-4 \citep{FIPS186_4}. Nó khác Schnorr ở chỗ \(r\) trong chữ
ký là \textbf{tọa độ \(x\) của điểm \(R\)} lấy modulo \(n\), thay vì bản
thân điểm \(R\).

\textbf{Ký:} Cho khóa riêng \(d\) và thông điệp \(m\):

\begin{enumerate}
\def\labelenumi{\arabic{enumi}.}
\tightlist
\item
  \(k \xleftarrow{R} [1, r-1]\)
\item
  \(R = [k]G\)
\item
  \(r_{\text{sig}} = R_x \pmod{n}\) (nếu \(r_{\text{sig}} = 0\), thử
  lại)
\item
  \(e = H(m)\)
\item
  \(s = k^{-1}(e + r_{\text{sig}} \cdot d) \pmod{n}\) (nếu \(s = 0\),
  thử lại)
\item
  Chữ ký: \(\sigma = (r_{\text{sig}}, s)\)
\end{enumerate}

\textbf{Xác minh:}

\begin{enumerate}
\def\labelenumi{\arabic{enumi}.}
\tightlist
\item
  Kiểm tra \(r_{\text{sig}}, s \in [1, n-1]\)
\item
  \(e = H(m)\), \(w = s^{-1} \pmod{n}\)
\item
  \(u_1 = ew\), \(u_2 = r_{\text{sig}} \cdot w\), cả hai tính mod \(n\)
\item
  \(P = [u_1]G + [u_2]Q\)
\item
  Hợp lệ khi \(P \neq \mathcal{O}\) và \(P_x \pmod{n} = r_{\text{sig}}\)
\end{enumerate}

\subsection{Cấu trúc dữ
liệu}\label{cux1ea5u-truxfac-dux1eef-liux1ec7u-1}

\begin{Shaded}
\begin{Highlighting}[]
\KeywordTok{pub} \KeywordTok{struct}\NormalTok{ EcdsaSignature }\OperatorTok{\{}
    \KeywordTok{pub}\NormalTok{ r}\OperatorTok{:}\NormalTok{ U1024}\OperatorTok{,}  \CommentTok{// r\_sig = R.x mod n, 128 byte}
    \KeywordTok{pub}\NormalTok{ s}\OperatorTok{:}\NormalTok{ U1024}\OperatorTok{,}  \CommentTok{// s in [1, n{-}1],     128 byte}
\OperatorTok{\}}
\end{Highlighting}
\end{Shaded}

\captionof{lstlisting}{Cấu trúc dữ liệu chữ ký ECDSA}\label{lst:ch4_ecdsa}\vspace{12pt}

Kích thước chữ ký: \(128 + 128 = \mathbf{256}\) byte --- nhỏ hơn Schnorr
vì không lưu điểm \(R\) đầy đủ.

\subsection{Cài đặt}\label{cuxe0i-ux111ux1eb7t-1}

Tương tự Schnorr, các thủ tục ký xác thực và kiểm chứng chuẩn mực của
thuật toán ECDSA cũng được triển khai trong không gian số học bằng các
thao tác đại số trực tiếp trên cấu trúc thành phần
\texttt{FieldElement}. Tại quy trình ký sinh, tiến trình tạo phần tử
ngẫu nhiên nonce được rà soát khắt khe và lập tức phải loại bỏ (chu kỳ
thử lại mới) nếu tổ hợp tạo ra vi phạm hệ số tiệm cận, cụ thể rơi vào 2
trường hợp suy biến tiềm ẩn: \(r_{\text{sig}} = 0\) hoặc \(s = 0\).
Logic bóc lỗi này tuân thủ trọn vẹn đặc tả bảo mật thuộc tiêu chuẩn
chuẩn hóa FIPS 186-4 được ban hành riêng cho ECDSA.

Chi tiết mã nguồn của các hàm \texttt{sign()} và \texttt{verify()} thuộc
thuật toán ECDSA được đính kèm tại mã nguồn công khai của dự án
\citep{LumenMath}.

\subsection{So sánh Schnorr và
ECDSA}\label{so-suxe1nh-schnorr-vuxe0-ecdsa}

\begin{longtable}[]{@{}
  >{\raggedright\arraybackslash}p{(\linewidth - 4\tabcolsep) * \real{0.3333}}
  >{\raggedright\arraybackslash}p{(\linewidth - 4\tabcolsep) * \real{0.3333}}
  >{\raggedright\arraybackslash}p{(\linewidth - 4\tabcolsep) * \real{0.3333}}@{}}
\caption{So sánh sơ đồ chữ ký Schnorr và ECDSA}\tabularnewline
\toprule\noalign{}
\begin{minipage}[b]{\linewidth}\raggedright
Tiêu chí
\end{minipage} & \begin{minipage}[b]{\linewidth}\raggedright
Schnorr
\end{minipage} & \begin{minipage}[b]{\linewidth}\raggedright
ECDSA
\end{minipage} \\
\midrule\noalign{}
\endfirsthead
\toprule\noalign{}
\begin{minipage}[b]{\linewidth}\raggedright
Tiêu chí
\end{minipage} & \begin{minipage}[b]{\linewidth}\raggedright
Schnorr
\end{minipage} & \begin{minipage}[b]{\linewidth}\raggedright
ECDSA
\end{minipage} \\
\midrule\noalign{}
\endhead
\bottomrule\noalign{}
\endlastfoot
Chuẩn hóa & ISO/IEC 14888-3 & FIPS 186-4, IEEE P1363 \\
Kích thước chữ ký & 384 byte (\(2 \times 128 + 128\)) & 256 byte
(\(2 \times 128\)) \\
Phép nhân vô hướng ký & 1 & 1 \\
Phép nhân vô hướng xác minh & 2 (tuần tự) & 2 (tuần tự) \\
Bảo mật chứng minh & RO Model, DL & Heuristic, DL + hash \\
Hỗ trợ tổng hợp & Có (MuSig, FROST) & Không \\
Rủi ro nonce tái sử dụng & Lộ khóa riêng & Lộ khóa riêng \\
\end{longtable}

\section{Chữ ký tổng hợp
BLS}\label{chux1eef-kuxfd-tux1ed5ng-hux1ee3p-bls}

\subsection{Lý thuyết (tổng
quan)}\label{luxfd-thuyux1ebft-tux1ed5ng-quan}

Chữ ký BLS (Boneh-Lynn-Shacham, 2001) \citep{BLS2001} là sơ đồ
\textbf{tổng hợp phi tương tác}: \(n\) người ký với \(n\) khóa riêng
khác nhau có thể tạo ra \textbf{một chữ ký duy nhất} đại diện cho tất
cả, và xác minh chỉ cần một phép ghép cặp thay vì \(n\) phép.

\textbf{Nền tảng:} Dùng phép ghép cặp song tuyến
\(e: \mathbb{G}_1 \times \mathbb{G}_2 \to \mathbb{G}_T\) với:

\begin{itemize}
\tightlist
\item
  \(\mathbb{G}_1 = E(\mathbb{F}_p)\),
  \(\mathbb{G}_2 = E'(\mathbb{F}_{p^2})\) (hoặc tương đương)
\item
  \(\mathbb{G}_T \subset \mathbb{F}_{p^k}^*\)
\end{itemize}

\textbf{Ký:} Người thứ \(i\) với khóa riêng \(d_i\):
\[\sigma_i = [d_i] H(m) \in \mathbb{G}_1\] trong đó
\(H: \{0,1\}^* \to \mathbb{G}_1\) là hàm băm lên nhóm điểm.

\textbf{Tổng hợp:} \(\sigma = \sigma_1 + \sigma_2 + \cdots + \sigma_n\)
(phép cộng điểm, không cần tương tác)

\textbf{Xác minh tổng hợp:}
\[e(\sigma, G_2) = e(H(m), Q_1 + Q_2 + \cdots + Q_n)\] trong đó
\(Q_i = [d_i]G_2\) là khóa công khai của người \(i\).

Tính đúng đắn:
\[e\!\left(\sum_i [d_i]H(m),\, G_2\right) = e\!\left(H(m),\, \sum_i [d_i]G_2\right) = e(H(m), \sum Q_i) \checkmark\]

\subsection{Vai trò của đường cong
KSS18}\label{vai-truxf2-cux1ee7a-ux111ux1b0ux1eddng-cong-kss18}

Đường cong được xây dựng trong luận văn này (họ Kachisa-Schaefer-Scott
với \(k = 18\)) \citep{CocksPinchK5K8} cung cấp nền tảng để cài đặt BLS:

\begin{itemize}
\tightlist
\item
  \textbf{\(\mathbb{G}_1 \subset E(\mathbb{F}_p)\)}: nhóm điểm trên
  trường cơ sở, kích thước \(r\).
\item
  \textbf{\(\mathbb{G}_T \subset \mathbb{F}_{p^{18}}^*\)}: nhóm đích của
  phép ghép cặp, bậc \(r\), kích thước trường 18432 bit --- đủ để kháng
  cả NFS lẫn TNFS.
\end{itemize}

Phép ghép cặp Optimal Ate (trên KSS18) chạy vòng lặp Miller với
\(O(\log r)\) bước, mỗi bước thực hiện phép tính trong
\(\mathbb{F}_{p^{18}}\). Cấu trúc NTT-friendly của \(p\)
(\(\text{two-adicity}(p-1) \geq 32\)) giúp tăng tốc các phép nhân trong
tháp trường.

\begin{quote}
\textbf{Lưu ý:} Cài đặt phép ghép cặp đầy đủ (vòng lặp Miller + phép lũy
thừa lần cuối) nằm ngoài phạm vi của luận văn này và là hướng phát triển
tiếp theo. Chương này tập trung vào Schnorr và ECDSA trên
\(\mathbb{G}_1\).
\end{quote}

\section{\texorpdfstring{Cài đặt hệ thống ký
(\texttt{curve1024-sig})}{Cài đặt hệ thống ký (curve1024-sig)}}\label{cuxe0i-ux111ux1eb7t-hux1ec7-thux1ed1ng-kuxfd-curve1024-sig}

Binary \texttt{curve1024-sig} cung cấp giao diện dòng lệnh kiểu GPG cho
toàn bộ vòng đời khóa:

\begin{Shaded}
\begin{Highlighting}[]
\ExtensionTok{curve1024{-}sig}\NormalTok{ keygen [{-}o }\OperatorTok{\textless{}}\NormalTok{file}\OperatorTok{\textgreater{}}\NormalTok{]}
\ExtensionTok{curve1024{-}sig}\NormalTok{ sign   }\AttributeTok{{-}f} \OperatorTok{\textless{}}\NormalTok{file}\OperatorTok{\textgreater{}}\NormalTok{ [{-}k }\OperatorTok{\textless{}}\NormalTok{key}\OperatorTok{\textgreater{}}\NormalTok{] [{-}{-}scheme schnorr}\KeywordTok{|}\ExtensionTok{ecdsa]}
\ExtensionTok{curve1024{-}sig}\NormalTok{ verify }\AttributeTok{{-}f} \OperatorTok{\textless{}}\NormalTok{file}\OperatorTok{\textgreater{}}\NormalTok{ [{-}k }\OperatorTok{\textless{}}\NormalTok{pub}\OperatorTok{\textgreater{}}\NormalTok{]}
\ExtensionTok{curve1024{-}sig}\NormalTok{ export{-}pub [{-}k }\OperatorTok{\textless{}}\NormalTok{priv}\OperatorTok{\textgreater{}}\NormalTok{] [{-}o }\OperatorTok{\textless{}}\NormalTok{pub}\OperatorTok{\textgreater{}}\NormalTok{]}
\end{Highlighting}
\end{Shaded}

\captionof{lstlisting}{Giao diện dòng lệnh CLI curve1024-sig}\label{lst:ch4_cli}\vspace{12pt}

Các tham số được đọc từ \texttt{config/curve1024.toml} thông qua
\texttt{build.rs} để sinh hằng số tĩnh --- đảm bảo zero-cost abstraction
và không có chi phí khởi tạo runtime.

\FloatBarrier
